% =====================================================================================================
\section{The posterior is right about the distribution}\label{sec:distribution}

Selecting the right model is one goal; predicting well is another. This section shows that the posterior recovers the
conditional law of the response given the predictors, at the rate $\rn$, in the distance used by \citet{SongLiang2015},
with no beta-min condition and no fixed lower bound on eigenvalues. The same holds for the part of the posterior that a
completed run returns.

\subsection{Setting for this section}\label{sec:dist-setting}

This section has its own setting, closer to that of \citet{SongLiang2015}: the observations are i.i.d.\ pairs, and the
number of observed predictors $\praw$ is used directly. In this section only, $x$ and $y$ denote one raw observation of
the predictors and of the response, $x_j$ is the $j$-th coordinate of $x$, and $(x^{(i)},y^{(i)})$ is the $i$-th observation; the
standardised data of Section~\ref{sec:standardise} do not appear.

The data are $n$ independent copies $(x^{(i)},y^{(i)})$, $i=1,\dots,n$, of a pair $(x,y)$: $x\in\R^{\praw}$ is the vector of all
observed predictors (a row of the raw design) and $y\in\R$ is the response. We write $\Dn$ for the data set. The number
$\praw$ may grow with $n$ and may be much larger than $n$, and the law of $(x,y)$ may change with $n$. As everywhere,
the method runs on the $p\le\praw$ predictors that are not constant in the sample, with the prior of
Section~\ref{sec:prior}. We write $\Pi(\cdot\mid\Dn)$ for the resulting exact posterior of $(S,\beta,\sigma^2)$, mapped
back to the original units: $S$ is the set of included predictors, $\beta\in\R^{\praw}$ is the coefficient vector (zero
outside $S$ and on removed predictors), and $\sigma^2$ is the noise variance. The posterior is not defined if fewer than two predictors remain or the
response is constant in the sample. Under the assumptions below such data sets are rare, and they are counted in the
exceptional sets of the theorems. Outside those sets, $p=\praw$.

\paragraph{What is estimated.} A posterior draw $(S,\beta,\sigma^2)$ defines a conditional law of $y$ given $x$,
\[
  \Normal\big(\bar y+(x-\bar x)^\top\beta,\ \sigma^2\big),
\]
with the sample-mean intercept ($\bar y$ and $\bar x$ are sample means). The law of $x$ is \emph{not} estimated. Let $f$
be the joint density of $(x,y)$ in which $x$ has its true law and $y$ given $x$ has the law above, and let $f^*$ be the
true joint density. The true model is $\Sz:=\{j:\beta^*_j\neq0\}$. For two laws $P$ and $Q$,
$H^2(P,Q):=1-\int\sqrt{dP\,dQ}$ is the squared Hellinger distance, and for
two densities $d_H(f,g):=(1-\int\sqrt{fg})^{1/2}$, so that $d_H^2(f,g)$ is $H^2$ of the corresponding laws. Because $f$ and $f^*$ share the law of $x$,
\begin{equation}\label{eq:dH}
  d_H^2(f,f^*)=\E_x\Big[H^2\Big(\Normal\big(\bar y+(x-\bar x)^\top\beta,\sigma^2\big),\
  \Normal\big(\alpha^*+x^\top\bstar,\sigs^2\big)\Big)\Big]:
\end{equation}
the squared Hellinger distance between the fitted and the true conditional laws of $y$ given $x$, averaged over a new $x$
independent of the data. This is the distance of \citet[Theorems~1--2]{SongLiang2015}, who use
$\int(\sqrt f-\sqrt{f^*})^2=2d_H^2$.

\paragraph{Further notation.} $\P^*$ is the law of the data. For the predictors, $\sd$, $\Var$ and $\Corr$ are population
standard deviations, variances and correlation matrices; $\zeta_j:=(x_j-\E x_j)/\sd(x_j)$ is the $j$-th predictor
standardised in the population, and $\zeta:=(\zeta_1,\dots,\zeta_{\praw})$; $x_T$ is the sub-vector of $x$ indexed by a
set $T$.

\paragraph{Constants, level and window.} The integer $s\ge0$ is an upper bound on the number of nonzero
coefficients (not necessarily the exact count). It and the constants $A\ge1$, $L\ge1$, $c_\psi>0$ and $\bar b>0$ are
fixed. We fix a \emph{level} $M\ge1$. The contraction radius below is $\sqrt{M\log n/n}$ and the exceptional
probability is $n^{-c_{\mathrm H}M}$, so a larger $M$ gives a larger radius and a smaller exceptional probability. The \emph{window} is
\[
  q_n:=\min\big\{\praw,\ \lceil C_qM\log n\rceil\big\},\qquad C_q:=2\big(s(A+2)+4\big)/\log2,
\]
the largest number of predictors that the conditions below constrain jointly; it is of order $\log n$ for a fixed level.

\subsection{The posterior contracts around the true conditional law}\label{sec:dist-results}

\inwords{Suppose the observations are i.i.d., the response follows a sparse Gaussian linear model, the predictors have
light tails, no small group of predictors is almost exactly collinear, and no single predictor carries an unbounded
signal. Then, except on a set of data sets of probability at most $n^{-c_{\mathrm H} M}$, the posterior gives probability at most
$n^{-c_{\mathrm H} M}$ to draws whose conditional law of $y$ given $x$ is farther than $\sqrt{M\log n/n}$ from the true one. That is,
the posterior contracts around the true conditional law at the rate $\rn$.}

\begin{localassumption}[Distribution-level setting]\label{asm:dist}
\begin{asmitems}
  \item\label{asm:dist-model} \emph{Gaussian linear model for i.i.d.\ pairs.} $y=\alpha^*+x^\top\bstar+\varepsilon$ with
    $\varepsilon\sim\Normal(0,\sigs^2)$, $\sigs^2>0$, independent of $x$. The numbers $\alpha^*$, $\bstar$, $\sigs$ and the
    law of $x$ may depend on $n$.
  \item\label{asm:dist-sparse} \emph{Sparsity and dimension.} At most $s$ coordinates of $\bstar$ are nonzero, and
    $2\le\praw\le n^A$.
  \item\label{asm:dist-subg} \emph{Sub-Gaussian predictors on sparse directions.} Every predictor has positive variance,
    and for every $u\in\R^{\praw}$ with at most $q_n$ nonzero entries and every $\lambda\in\R$,
    $\E\exp(\lambda u^\top\zeta)\le\exp\big(L^2\lambda^2\Var(u^\top\zeta)/2\big)$.
  \item\label{asm:dist-nec} \emph{No near-exact collinearity.} $\lmin(\Corr(x_T))\ge c_\psi/n$ for every set $T$ of at
    most $q_n$ predictors.
  \item\label{asm:dist-bs} \emph{Bounded signal.} $\abs{\beta^*_j}\sd(x_j)\le\bar b\,\sigs$ for every $j$.
\end{asmitems}
\rulesout{(a) non-Gaussian or heteroscedastic noise and dependent observations; (b) dense coefficient vectors and more
than polynomially many predictors; (c) heavy-tailed predictors; (d) groups of $O(\log n)$ predictors whose
correlation matrix has smallest eigenvalue below $c_\psi/n$, while predictors may otherwise be very strongly correlated; (e) a single predictor whose signal is
unboundedly large compared with the noise, and huge coefficients that cancel.}
\end{localassumption}

No lower bound on the nonzero coefficients (no beta-min condition) and no fixed lower bound on eigenvalues is assumed.

\begin{theorem}[Distribution-level consistency]\label{thm:dist}
There are constants $c_{\mathrm H}>0$, $M_0$ and $n_0$, depending only on $(s,A,L,c_\psi,\bar b)$, such that the following
holds. Fix a level $M\ge M_0$ and assume Assumption~\ref{asm:dist}, with the window $q_n$ of this $M$. Then for every
$n\ge n_0$ with $\sqrt n/\log n\ge M$,
\begin{equation}\label{eq:dist-i}
  \P^*\Big[\,\Pi\Big(d_H(f,f^*)>\sqrt{M\log n/n}\ \Big|\ \Dn\Big)>n^{-c_{\mathrm H} M}\Big]\ \le\ n^{-c_{\mathrm H} M}.
\end{equation}
\end{theorem}

\status{\textbf{Proved and independently checked}. Verbatim: Appendix~\ref{app:status}, row 2.1.}

\inwords{The same holds for every level $M$ up to $\sqrt n/\log n$, with the same constants, and then the posterior also
puts almost no mass on models with more than $s+C_2M$ predictors. Without the collinearity condition, the radius grows by
a factor $\sqrt{\log(\praw\vee n)}$.}

\begin{proposition}[Level-$M$ form, and the case without the collinearity condition]\label{prop:dist}
\textup{(i)} There is a constant $C_2$, depending only on $(s,A,L,c_\psi,\bar b)$, such that the following holds. The
constants of Theorem~\ref{thm:dist} serve every level $M$ with $M_0\le M\le\sqrt n/\log n$, and $M$ may now depend on $n$
(Theorem~\ref{thm:dist} fixes it): whenever $n\ge n_0$ and Assumption~\ref{asm:dist} holds with the window $q_n$ of this $M$,
\eqref{eq:dist-i} holds, and also
\begin{equation}\label{eq:dist-ii}
  \P^*\Big[\,\Pi\big(\abs{S}>s+C_2M\ \big|\ \Dn\big)>n^{-c_{\mathrm H} M}\Big]\ \le\ n^{-c_{\mathrm H} M}.
\end{equation}
The two exceptional sets of data sets can be taken to be the same set, and the bounds are uniform over all laws of the
data that satisfy Assumption~\ref{asm:dist} with the given constants.

\textup{(ii)} If Assumption~\ref{asm:dist}\ref{asm:dist-nec} is dropped and Assumption~\ref{asm:dist}\ref{asm:dist-subg} is
required only for single predictors ($u$ with one nonzero entry), then
\begin{equation}\label{eq:dist-iii}
  \P^*\Big[\,\Pi\Big(d_H(f,f^*)>\sqrt{M\log n\cdot\log(\praw\vee n)/n}\ \Big|\ \Dn\Big)>n^{-c_{\mathrm H} M}\Big]\ \le\ n^{-c_{\mathrm H} M},
\end{equation}
and~\eqref{eq:dist-ii} holds with $s+C_2M\log n$ in place of $s+C_2M$.
\end{proposition}

\status{\textbf{Proved and independently checked} (part of the same result as Theorem~\ref{thm:dist}). Verbatim:
Appendix~\ref{app:status}, row 2.1.}

\begin{remark}[On the assumptions]\label{rem:dist-asm}
(i) \emph{Designs that satisfy the assumptions include} jointly Gaussian predictors, which satisfy
Assumption~\ref{asm:dist}\ref{asm:dist-subg} with $L=1$; independent predictors, which satisfy it when each does with a
common constant $L$; and, for Assumption~\ref{asm:dist}\ref{asm:dist-nec} for all large $n$, every correlation matrix
whose $q_n\times q_n$ principal blocks have smallest eigenvalue at least a fixed $\psi>0$, for instance correlation
$\varrho^{\abs{j-k}}$ with $\abs{\varrho}<1$, or equal correlation $\varrho\in[0,1)$.
(ii) Assumption~\ref{asm:dist}\ref{asm:dist-subg} is used only to make the sample variances and covariances of small groups of
predictors track their population values with probability $1-n^{-c_{\mathrm H} M}$; the passage from the sample to the population
distance uses no tail condition.
(iii) Assumption~\ref{asm:dist}\ref{asm:dist-nec} keeps in force the penalty that the prior on the coefficients imposes on
each extra predictor (the Occam factor, of order $n^{-1/2}$); this gives the rate $\rn$, the level $n^{-c_{\mathrm H} M}$ and the bounded model size. Whether it is needed for the rate
$\rn$ is open.
(iv) Assumption~\ref{asm:dist}\ref{asm:dist-bs} bounds the signal in units of the noise. The noise level $\sigs$ is free because
rescaling the response leaves the posterior over models unchanged and rescales the posterior of $(\beta,\sigma^2)$.
\end{remark}

\subsection{The parametric rate, and the output of a run}\label{sec:dist-param}

The logarithm in the rate disappears when the true predictors are detectable.

\inwords{If, in addition, small groups of predictors have a fixed eigenvalue floor and every true signal is above the
detection scale $\rn$, the posterior contracts at the parametric rate $1/\sqrt n$, up to any factor that tends to
infinity.}

\begin{localassumption}[Detectable signal]\label{asm:dist-det}
(a) $\lmin(\Corr(x_T))\ge\psi$ for a fixed $\psi>0$ and every set $T$ of at most $(2A+3)s+2$ predictors; and (b)
$\min\{\abs{\beta^*_j}\sd(x_j)/\sigs:\beta^*_j\neq0\}\ge b_n$, where $b_n>0$ and $b_n^2n/\log n\to\infty$.
\rulesout{collinearity among small sets of predictors, and true signals below the detection scale $\rn$.}
\end{localassumption}

\begin{theorem}[Parametric rate when the true predictors are detectable]\label{thm:dist-param}
Assume Assumption~\ref{asm:dist} with the window $q_n$ of level $M=1$, with exactly $s\ge1$ nonzero coordinates of $\bstar$,
and Assumption~\ref{asm:dist-det}. Then for every sequence $M_n\to\infty$,
\[
  \Pi\big(d_H(f,f^*)>M_n/\sqrt n\ \big|\ \Dn\big)\to0\quad\text{in $\P^*$-probability.}
\]
\end{theorem}

\status{\textbf{Proved and checked, conditional on} Theorem~\ref{thm:post-fixed} (in-probability form), which is
proved and independently verified (two independent verification passes). Verbatim: Appendix~\ref{app:status}, row 2.1.}

\paragraph{What a completed run returns.} The method explores the models with a sequential Monte Carlo sampler; it does
not compute the full posterior. (Small problems that it solves by exact enumeration are not covered here.) Runs that stop
early are not covered either. The output of a completed run contains a list $\mathcal T$ of at most $40p+1$ distinct
models (a list of size $O(p)$; the constant is the code's). The list holds the best-scoring models among a pool: the
models scored in the final sampler run, plus those carried into it, namely the precomputed block tables and the models
kept from the calibration. Each model in the list comes with its log score $F(S)$, the logarithm of its unnormalised
posterior weight. The output
contains no draws of $(\beta,\sigma^2)$. Let $\pi(\mathcal T\mid\Dn)$ be the posterior probability of the list, and let
$\hat\Pi$ (not computed by the method) be the posterior restricted to the list: the models in $\mathcal T$, weighted in
proportion to the exponentials of their returned log-weights, each with its exact conditional posterior of
$(\beta,\sigma^2)$.

\inwords{If a completed run's list carries a non-negligible share of the posterior, the posterior restricted to that list
is as good as the full posterior: it, too, puts almost no mass outside the ball of radius $\sqrt{M\log n/n}$.}

\begin{corollary}[The output of the sampler]\label{cor:dist-output}
Assume the conditions of Proposition~\ref{prop:dist}(i), and that the returned log-weights are computed in exact arithmetic.
With $\P^*$-probability at least $1-n^{-c_{\mathrm H} M}$ over the data, every completed run, whatever its random numbers, with
$\pi(\mathcal T\mid\Dn)\ge n^{-c_{\mathrm H} M/2}$ satisfies
\[
  \hat\Pi\Big(d_H(f,f^*)>\sqrt{M\log n/n}\Big)\ \le\ n^{-c_{\mathrm H} M/2}.
\]
If a completed run scored the true model $\Sz$ and $\pi(\Sz\mid\Dn)>\tfrac12$, then the list contains $\Sz$, the
reported model is $\Shat=\Sz$, and $\pi(\mathcal T\mid\Dn)>\tfrac12$.
\end{corollary}

\status{\textbf{Proved and checked, under exact arithmetic}. The description of the output was checked against the code
by independent reviewers. Verbatim: Appendix~\ref{app:status}, row 2.1.}

\begin{remark}[What is not claimed]\label{rem:dist-output}
(i) Whether the list carries enough posterior probability is a property of the run: Corollary~\ref{cor:dist-output} does
not assert it, and the output does not certify it.
(ii) Theorem~\ref{thm:dist} and Corollary~\ref{cor:dist-output} use no result from other sections (the corollary concerns
the exact posterior restricted to the returned list);
Theorem~\ref{thm:dist-param} also uses the posterior concentration theorem (Theorem~\ref{thm:post-fixed}). Statements that
a completed run has in fact found and kept the true model rest, in addition, on the sampler results of
Section~\ref{sec:delivery} and on exact arithmetic.
(iii) Statement~\eqref{eq:dist-ii} bounds the size of the selected models, not their identity.
\end{remark}

\subsection{Fixed designs and the in-sample distance}\label{sec:dist-fixed}

Theorem~\ref{thm:dist} is about i.i.d.\ observations and does not apply to a fixed design. For a fixed design,
Remark~\ref{rem:dist-fixed} gives contraction in an in-sample distance instead. In this remark we return to the notation
of Section~\ref{sec:setting}.

\begin{remark}[Further consequences: fixed designs]\label{rem:dist-fixed}
Let the design be fixed, and let $\bar h_n$ be the in-sample Hellinger distance: $\bar h_n^2$ is the average, over the $n$
observations, of the squared Hellinger distance between the fitted and the true laws of one observation, both centred over
the sample (so $\bar h_n$ is a root mean square). Assume \as{GL}, $\sz\ge1$ fixed and
$2\le p\le n^A$, the energy bound $\norm{\mu}^2\le nM_\mu$, and the floor $\lmin(\Rhat_{\Sz})\ge\varphi_0^2$ on the true
model. All limits below are in probability. (a) If also $1-\pis{\Sz}\to0$, then $\Pi(\bar h_n>M_n/\sqrt n\mid y)\to0$ for
every $M_n\to\infty$: the distance contracts at rate $M_n/\sqrt n$. (b) Without beta-min, under the superset Occam-sum
condition of Remark~\ref{rem:post-fixed-asm}(iii) instead (which \as{NC} implies for all large $n$),
$\Pi(\bar h_n^2>C_H\log n/n\mid y)\to0$ for a constant $C_H$: the squared distance contracts at rate $\log n/n$.
\end{remark}

\status{\textbf{Proved and checked}; its inputs are Theorem~\ref{thm:post-fixed} in probability form (proved and
independently verified) and the random-design transfer lemma (proved and independently checked). Verbatim: Appendix~\ref{app:status}, row 2.2.}

\subsection{Comparison with Song and Liang (2015), Theorems 1--2}\label{sec:SL-dist}

Page and equation numbers are those of the published paper \citep{SongLiang2015}. In this subsection the symbols $r_n$,
$\bar r_n$, $P_n$, $\xi$, $\alpha$, $\theta$, $k$ and the constant $c_1$ are Song and Liang's, and their displays and
conditions are numbered as in their paper; $b_n\succ a_n$ means $a_n/b_n\to0$.

\paragraph{The analogue is display~\eqref{eq:dist-i}.} Their distance is
$d(f_1,f_2)=\{\int\!\!\int(f_1^{1/2}-f_2^{1/2})^2\,dx\,dy\}^{1/2}$ between joint densities of $(x,y)$ that share the law of
$x$ (\S2.1.2, p.~952), so $d^2=2d_H^2$: the same distance. Their display~(15) (Theorem~1, p.~952; Theorem~2 states the
same form, p.~953) reads $P^*[\pi\{d(f,f^*)>\varepsilon_n\mid D_n\}>\exp(-c_1n\varepsilon_n^2)]<\exp(-c_1n\varepsilon_n^2)$
(their constant $c_1$); display~\eqref{eq:dist-i} has this form with $n\varepsilon_n^2=M\log n$ (and $d_H$ in place of
$d$).

\paragraph{Stronger.} (1) The prior is the one the method uses, with no tuning. Their theorems assume independent
Bernoulli$(r_n/P_n)$ inclusion with a hard cap $\abs{\xi}\le\bar r_n$ on the model size and conditions (9)--(14)
(pp.~951--953), which the method's beta-binomial prior with cap $\min(p,\lfloor n/2\rfloor)$ does not meet. (2) The radius
is a constant times $\rn$. At polynomial $p$ their Theorem~2 (p.~953) needs $\bar r_n\log n=o(n\varepsilon_n^2)$ with
$\bar r_n\ge r_n\to\infty$ (their $r_n$) (and their (14)), so $n\varepsilon_n^2/\log n\to\infty$; their Corollary~1 (p.~954) gives
$\varepsilon_n=O\{n^{-(1-\alpha)/2}\log^{k/2}n\}$. In display~\eqref{eq:dist-i} the radius constant does not depend on $s$,
and the smallest level $M_0$ grows at most linearly in $s$ (in the general form of the theorem, where the in-sample
signal energy and coefficient cost are bounded by fixed constants; under Assumption~\ref{asm:dist}\ref{asm:dist-bs} those bounds themselves grow
with $s$). (3) Theorem~\ref{thm:dist-param} gives the parametric rate $M_n/\sqrt n$ and Corollary~\ref{cor:dist-output}
covers the sampler's output; their results concern the exact posterior only and give no parametric rate. (4) No lower
bound on the growth of $p$ is needed; they assume $P_n\succ n^\theta$ (p.~950).

\paragraph{Weaker.} (1) $\praw\le n^A$; their Corollary~1 allows $P_n<\exp(Cn^\alpha)$. (2) Exact sparsity; their
Theorem~1 covers approximate sparsity (their (3): $\lim_n\sum_j\abs{\beta^*_j}<\infty$, with predictors bounded by 1), under
which the number of nonzero coefficients may grow. (3) $s$ is fixed, with constants depending on $s$; their Theorem~2 also
has a true model of bounded size (their (4), p.~950), so this is weaker only against their Theorem~1. (4)
Assumption~\ref{asm:dist}\ref{asm:dist-nec} is a condition on the design that their Theorems~1--2 do not have (their eigenvalue
conditions, (11) and (14), are on the prior covariance).

\paragraph{Not the analogue: display~\eqref{eq:dist-iii}.} It is the fallback without Assumption~\ref{asm:dist}\ref{asm:dist-nec}. Its radius is larger
by $\sqrt{\log(\praw\vee n)}$, so its exponent is $c_{\mathrm H} n\varepsilon_n^2/\log(\praw\vee n)$, not $c_{\mathrm H} n\varepsilon_n^2$; and, as
proved, in the general form of the theorem (in-sample signal energy and coefficient cost bounded by fixed constants), its
squared radius at the smallest admissible level is at most of order $(1+s)^3\log n\log(\praw\vee n)/n$ for fixed $s$ (an
upper bound; the proof is not known to be tight in $s$); under Assumption~\ref{asm:dist}\ref{asm:dist-bs} those bounds themselves grow with $s$.

\paragraph{Different, not ranked.} Tail conditions on the predictors: Assumption~\ref{asm:dist}\ref{asm:dist-subg} and Assumption~\ref{asm:dist}\ref{asm:dist-bs} here;
theirs are $\abs{x_i}\le1$ (Theorem~1), Condition~1 (an exponential-moment bound on averages of at most $\bar r_n$
predictors, Theorem~2, p.~953) and Condition~2 (second moments, for their limit form (16), p.~954).

\paragraph{Overall: incomparable.} Stronger on the prior, the rate and the coded output; weaker on the range of $p$, on
approximate sparsity and, for display~\eqref{eq:dist-i}, by Assumption~\ref{asm:dist}\ref{asm:dist-nec}.
